

\section{The Rees-Sushkewitch Theorem}
\label{sec:2.3}

\begin{lemma}  	\label{lem:2.1}
Let $S$ be a simple semigroup.

     If $e \in S$ is a primitive idempotent, then $Se,\ eS$ are
minimal left and right ideals respectively and the set $eSe$ is a
group.  In addition, $eSe = eS \cap Se$.
\end{lemma}

\begin{proof}  
Consider $eS$.  
Clearly it is a right ideal.
Suppose $I \subset eS$ is a right ideal.  
Let $a \in I$.

-- Notice that for $a \in eS$, $e$ acts like left identity.
\[
     a \in eS \Rightarrow a = es_1 \Rightarrow ea
          = ea = e^2s_1 = es_1 = a.
\]
Now $SaS = S$ since $S$ is simple.  Thus $xay = e$ for some $x,y \in
S$.  Notice $(exe)(ea)(ye) = e$.  Let $exe = x'$, $ye = y'$, and
$z' = ay'x'$.  Then
\[
     \aligned
     z'z' &= ay'x' ay'x'                                \\
     &= ay'(exe)(ea)(ye)x'                              \\
     &= ay'ex'                                          \\
     &= ay'x' = z'.
     \endaligned
\]
Thus $z'$ is an idempotent.  
Also
\[
     \gathered
     z'e = (ay'x')e = ay'x' = z'                        \\
     ez' = eay'x' = ay'x' = z'.
     \endgathered
\]
But $e$ was assumed to be primitive, thus $e = z'$.  
But $e = ay'x' \in aS \subset IS \subset I$, thus
\[
     \gathered
     eS \subset IS \subset I                            \\
     \Rightarrow eS = I
     \endgathered
\]
and $eS$ is minimal.
\qed
\end{proof}

\runinhead{Homework.}  
Show that $Se$ is a minimal left ideal.

     Now let $G = eSe$, and $g \in G$.
Then
\[
     \aligned
     G g = G(ese) &= (eSe)(ese)                    \\
     &= e(Sese)                                         \\
     &= eSe = G \quad \hbox{by minimality}              \\
     gG = (ese)G &= (eseS)e = eSe = G.
     \endaligned
\]
Thus $G$ is a group, since it is left simple and right simple.

\runinhead{Homework.} Show that $eSe = eS \cap Se$.

\medskip
     In fact, we have the following.

\begin{theorem}  	\label{thm:2.1}
Let $S$ be a finite semigroup, with $R$ a minimal right ideal and $L$ a minimal left ideal.  
Then $RL$ is a group.
\end{theorem}

\begin{proof}  
Let $RL = G$, $g \in G$.  
Then
\[
     \aligned
     gG &= rlRL                                         \\
     &= RL
     \endaligned
\]
since $(rl)R \subset R$ and is a right ideal.  Similarly,
\[
     \aligned
     Gg &= RLrl \qquad rl \in L                         \\
     &= RL.
     \endaligned
\]
Thus $RL$ is both left and right simple and thus is a group.
\qed
\end{proof}

\begin{lemma}    	\label{lem:2.2}
Let $L$ be a minimal left ideal in a semigroup $S$.  
Then $Lc$ is a minimal left ideal for all $c \in S$.
\end{lemma}

\begin{proof}  
Suppose not.  
Let $A \subsetneq Lc$ be a left ideal.  
Let $B = \{ x \in L : xc \in A\} \subset L$.  
By definition, $Bc \subset A \subsetneq Lc$.  
Notice that if $a \in A$, $a = lc$ for some $l \in L \Rightarrow l \in B \Rightarrow lc = a \in Bc \Rightarrow Bc = A$.

Now let $b \in B$, $s \in S$.  
Then,
\[
     (sb)c \subset sA \subset A
               \qquad (A \hbox{ is a left ideal)}
\]
Also,
\[
     sb \subset sL \subset L.
\]
Thus
\[
     \alignedat{2}
     &(sb)c \in A &\qquad &\hbox{where } sb \in L     \\
     \Rightarrow \quad
          &(sb) \in B &\qquad &\hbox{for all } s \in S     \\
     \Rightarrow \quad
          &B &\qquad &\hbox{is a left ideal}              \\
     \Rightarrow \quad
          &B = L &\qquad &\hbox{by minimality}            \\
     \Rightarrow \quad
          &Bc = Lc = A. &\quad & \qquad\otimes
     \endalignedat
\]
Thus
\[
     Lc \hbox{ is minimal}. \quad \qed
\]
\end{proof}

\runinhead{Homework.}  
Prove the analogous result for right ideals

\begin{lemma}    	\label{lem:2.3}
Let $S$ be completely simple.  
Then $S = \bigcup\limits_i L_i = \bigcup\limits_j R_j$ where $\{ L_i\}$,
$\{ R_j\}$ are the set of all minimal left and right ideals,
respectively.
\end{lemma}

\begin{proof}  
Notice that the union of two left ideals is a left
ideal.  (Homework)
\end{proof}

Let $A = \bigcup\limits_i L_i$ where each $L_i$ is a minimal left ideal.  We will show $A$ is a right ideal as well.  
Now
\[
     \aligned
     &a \in L_i \qquad \hbox{for some } i               \\
     \hbox{and} \qquad
          &as \in L_is \subset A,
     \endaligned
\]
since $L_i s$ is a minimal left ideal by Lemma 2.  
Thus
\[
     as \in A \quad \hbox{for all} \quad s \in S.
\]
Thus $A$ is a right ideal, therefore an ideal.  But $S$ is simple
$\Rightarrow A = S$.

\runinhead{Notation.}  
For any subset $T \subset S$, $E(T)$ will be the set of idempotents of $T$.  
Consider $E(L)$ where $L$ is a minimal left ideal and let $e \in E(L)$.  
Then since $e \in L$,
\[
     Se \subset SL \subset L
\]
and $Se$ is a left ideal $\Rightarrow Se = L$.  
(In fact, every minimal left ideal is of the form $Ss$ for some $s$.)

Now let $l \in L$.  
Then $l = se$ for some $s$ and $le = (se)e = se = l$.  
Thus {\it each} idempotent of $L$ acts like a right identity of $L$.  Take $e_1,e_2 \in E(L)$.  
Suppose
\[
     \gathered
     e_1e_2 = e_2e_1                                    \\
     \Rightarrow e_1 = e_2.
     \endgathered
\]
Thus {\it all} idempotents are primitive.  (Since $S = \bigcup L_i$,
every idempotent is in some $L$.)

\runinhead{Note.}  
$E(L)$ is a left zero semigroup.  
Similarly, $E(R)$ is a right zero semigroup.

\runinhead{Rees-Sushkewitch Theorem.}  
A semigroup is completely simple
iff there exist sets $X,$ $Y$, and a group $G$ such that $S \cong X \times G \times Y$, where $X \times G \times Y$ is a Rees product semigroup.

\begin{proof}  Let $S$ be completely simple and $e \in E(s)$.
Let $X = E(Se)$, $Y = E(eS)$, $G = eSe$.  
Recall we have shown that $G$ is a  group.  
Then from the previous discussion $X = E(Se)$.
$Y = E(eS)$ are left zero and right zero semigroups respectively.

     Define $\psi: X \times G \times Y \to S \qquad (x,g,y) \mapsto xgy$.
Suppose that $\psi(x_1,g_1,y_1) = \psi(x_2,g_2,y_2)$.
\[
     \alignedat{2}
     \hbox{Then}& &\quad
          x_1g_1y_1 &= x_2g_2y_2                        \\
     \hbox{and}& &\quad
          ex_1g_1y_1e &= ex_2g_2y_2e                    \\
     y_ie = e, \ ex_i = e& &
          eg_1e &= eg_2e                                \\
     \Rightarrow& &\quad
          g_1 &= g_2.
     \endalignedat
\]
Also
\[
     \alignedat{2}
     & &\quad x_1gy_1 &= x_2gy_2                        \\
     &\Rightarrow &\quad
          x_1gy_1eg^{-1} &= x_2gy_2eg^{-1}              \\
     &\Rightarrow &\quad
          x_1e &= x_2e                                  \\
     &\Rightarrow &\quad
          x_1 &= x_2.
     \endalignedat
\]
Similarly, $y_1 = y_2$.  
Thus, $\psi$ is one-to-one.

     Now let $s \in S$.  
Associate with $s$ the following element in $X \times G \times Y$
\[
     \aligned
     \a: &s \to X \times G \times Y                             \\
     &s \mapsto \bigl[ E(sSe),\, ese,\, E(eSs) \bigr]   \\
     &\qquad = (e_1,\, ese,\, e_2)
     \endaligned
\]
(well defined since $E(sSe)$, $E(eSs)$ are unique).  Notice that
$[(s)\a]\psi$ is the identity since
\[
     \aligned
     e_1(ese)e_2 &= e_1se_2 = s                         \\
     \hbox{since} \quad e_1 &\in E(sS)
          \qquad \hbox{and acts like a left identity}   \\
     e_2 &\in E(Ss) \qquad
          \hbox{and acts like a left identity}.
     \endaligned
\]
Thus, $\psi$ is one-to-one and onto.  
\qed
\end{proof}

Now, let $S$ be any finite semigroup and let
$I = \{ I_1,I_2,\ldots,I_n\}$ be the set of all ideals of $S$.  
Let $K = \bigcap\limits^n_{j=1} I_j$.  
Then $K$ is the unique minimal ideal of $S$.  
By construction $K$ is a simple semigroup.  
To show $K$ is completely simple we must show it contains a primitive idempotent.
Now $K$ is also a semigroup, and in particular is a left ideal, and a
right ideal.  
Since there are only a finite number of left ideals possible, if we order them by inclusion there must exist minimal elements.  
Similarly, minimal right ideals exist.  

%\begin{lemma}    	\label{lem:2.4.} 
%Let $R$ be a minimal right ideal and $L$ a minimal left ideal.  
%Then $RL$ is a group and its idempotent is primitive.
%\end{lemma}

%\begin{proof}  
%Let $G = RL$.  Then $G \cdot G = RLRL \subset RL = G$.  
%Thus $G$ is closed.
%
%Notice that $rlR = R$ by minimality.  
%Similarly, $Lrl = L$.  
%Thus if $g = rl$
%\[
%     \aligned
%     Gg &= (RL)(rl) = RL = G                            \\
%     \hbox{and} \qquad
%          gG &= (rl)RL = RL = G.
%     \endaligned
%\]
%Thus $G$ is a group.
%\qed
%\end{proof}

\runinhead{Homework.}  
Show that $RL = R \cap L$.  
Let $e$ be the identity of $RL$.  
Suppose $f = f^2$ is such that $ef = fe = f \Rightarrow \F \in R \cap L \Rightarrow e = f$ by uniqueness.

Combining the lemma with the previous discussion we have:

\begin{theorem}  	\label{thm:2.2}
Let $S$ be a finite semigroup.  
Then $S$ has a minimal ideal $K$ and $K$ is a completely simple semigroup.  
$K$ is often referred to as the kernel of $S$.
\end{theorem}
